\chapter{Microeconomics for Managers}\label{ch:microeconomics}

% Scope: elasticity, marginal analysis, market structure, and auction theory.
% Cross-link: elasticity <-> bid-response & MMM saturation (\cref{ch:paid-acquisition-platforms},
%   \cref{ch:attribution-and-measurement}); auction theory <-> the ad auction (\cref{ch:paid-acquisition-platforms}).

Markets allocate resources through prices; managers move prices without understanding how price
and quantity interact at the margin. Elasticity quantifies that interaction; market-structure
analysis locates the firm on the spectrum of competitive pressure; auction theory explains the
mechanism by which two of the largest media markets in history clear in milliseconds. Together
they give every pricing, positioning, and acquisition decision a microeconomic backbone.

\section{Price Elasticity and Marginal Analysis}\label{sec:elasticity}
% complexity: 3

If you raise the price, do you earn more or less total revenue? The answer depends entirely on
how sensitive demand is to price --- a question elasticity answers before you commit to the change.
Demand curves slope downward: higher price, fewer units. But the revenue effect of a price change
depends on which effect dominates --- the price-per-unit increase or the volume loss. A price rise
earns more total revenue only if customers are relatively insensitive (inelastic demand). In
elastic markets, the volume loss swamps the margin gain. Marginal analysis sharpens this: the
profit-maximising rule is to produce until the last unit's marginal revenue equals its marginal
cost --- any departure from that leaves money on the table.

Price elasticity of demand is the percentage change in quantity demanded per one-percent change in price:
\begin{equation}\label{eq:elasticity}
  \varepsilon \;=\; \frac{\partial Q/Q}{\partial P/P}
              \;=\; \frac{\partial Q}{\partial P}\cdot\frac{P}{Q}.
\end{equation}
\(\varepsilon<0\) always (law of demand). The revenue-maximising point is \(\varepsilon=-1\)
(unit-elastic); below it (inelastic, \(|\varepsilon|<1\)) a price rise raises revenue; above it
(elastic, \(|\varepsilon|>1\)) a price rise cuts revenue. Marginal revenue is linked to elasticity by:
\[
  \mathrm{MR} \;=\; P\!\left(1+\frac{1}{\varepsilon}\right).
\]
Setting \(\mathrm{MR}=\mathrm{MC}\) yields the standard profit-maximum. Revenue changes as:
\[
  \frac{\Delta R}{\Delta P} \;=\; Q\!\left(1 + \frac{1}{\varepsilon}\right),
\]
so raising price in an elastic market (\(|\varepsilon|>1\)) generates a negative revenue change ---
the volume loss dominates.

\Cref{eq:elasticity} assumes a static demand curve, constant tastes, and no strategic response
from competitors. Elasticity varies along the demand curve (it is not constant for a linear demand
function), so a calibration at one price point does not transfer to a very different price.
Long-run elasticities are generally higher than short-run ones: customers need time to find
substitutes, change habits, or switch platforms. In digital markets, platform lock-in and
switching costs can compress the measured elasticity far below the structural long-run value.

\begin{example}
The running SaaS business prices its base plan at \money{50}/month (\money{600}/year).
Suppose a demand study estimates \(\varepsilon=-1.5\) at the current price --- moderately elastic.
A proposed \qty{10}{\percent} price increase to \money{55}/month implies a volume decrease of
\(1.5\times10\%=\qty{15}{\percent}\):
\[
  \Delta R \;\approx\; Q\!\left(1+\frac{1}{-1.5}\right)\times\Delta P
           = Q\times\frac{1}{3}\times\money{5}.
\]
Revenue per existing customer rises by \money{5}, but \qty{15}{\percent} of customers churn ---
for a customer base of \num{200}, that is \num{30} lost subscriptions at \(\mathrm{LTV}=\money{3150}\)
each, a foregone value of \money{94500}. The price rise is unlikely to be net-positive unless the
elasticity estimate is materially wrong or a cohort of customers is genuinely inelastic. Run the
price experiment on a small segment before rolling out, and measure churn response directly
rather than relying on the aggregate \(\varepsilon\) estimate.
\end{example}

\begin{admonition}{Warning}
Using aggregate elasticity for a segment decision hides dangerous heterogeneity. An overall
\(\varepsilon=-1.5\) masks a distribution: some customers are nearly inelastic (enterprise, high
switching cost), others highly elastic (SMB trial users). A blanket price rise extracts from
neither group optimally. Segment-level pricing --- the subject of \cref{ch:pricing} --- is the
correct response; elasticity segmentation is its empirical foundation.
\end{admonition}

In media-mix modelling, the saturation curve --- diminishing returns to additional spend --- is
precisely a demand-side elasticity: the \emph{response elasticity} of conversions to media
dollars, measured at the current spend level. As spend increases, the marginal response declines,
which is why the \(\mathrm{iROAS}\) (\cref{eq:iroas}) estimated by \textcite{hanssens2001marketing}
and the MMM literature is a local elasticity estimate around the current budget point. The
adstock-saturation model in \cref{ch:attribution-and-measurement} is the time-distributed version
of \cref{eq:elasticity}: the same trade-off between spend and diminishing returns, extended over
multiple periods via the adstock decay (\cref{eq:adstock}).

Elasticity feeds directly into price architecture (\cref{ch:pricing}), campaign budget
optimisation (\cref{ch:paid-acquisition-platforms}), and the saturation curves of marketing-mix
modelling (\cref{ch:attribution-and-measurement}). \textcite{pindyck2017} provide the full
theoretical treatment; \textcite{mankiw2020} the undergraduate-level derivation.

\section{Market Structure and the Lerner Index}\label{sec:lerner}
% complexity: 3

How much pricing power does the business actually have? Market structure --- the number and
symmetry of competitors --- determines the upper bound on sustainable margin. The Lerner index
converts that structural position into a number. Perfect competition forces price to marginal
cost; pure monopoly allows maximum mark-up. Between those extremes lie oligopoly (few firms,
strategic interdependence) and monopolistic competition (many firms, differentiated products).
Most technology and SaaS businesses operate in monopolistic competition: pricing power exists
but is constrained by the threat of substitution. The Lerner index measures the fraction of
price that is mark-up, linking directly to elasticity: high pricing power requires inelastic
demand.

The Lerner index is the price--cost margin expressed as a fraction of price:
\begin{equation}\label{eq:lerner}
  L \;=\; \frac{P - \mathrm{MC}}{P} \;=\; -\frac{1}{\varepsilon}.
\end{equation}
The second equality follows from setting \(\mathrm{MR}=\mathrm{MC}\) at the profit maximum and
substituting \(\mathrm{MR}=P(1+1/\varepsilon)\). \(L\in[0,1]\): zero in perfect competition
(\(P=\mathrm{MC}\)), one in pure monopoly (\(\mathrm{MC}=0\)). The formula encodes the same
trade-off as \cref{eq:elasticity}: lower elasticity (less price sensitivity) permits a higher
Lerner index.

Market structure typology:
\begin{description}
  \item[Perfect competition] \(P=\mathrm{MC}\), \(L=0\). Many identical firms, zero economic profit in
    the long run. Commodity agriculture, some spot-market financial instruments.
  \item[Monopolistic competition] Differentiated products, free entry, \(L>0\) in the short run but
    eroded by imitation. Most consumer goods, SaaS with many close substitutes.
  \item[Oligopoly] Few large firms, strategic interdependence (each firm's price affects rivals).
    Telecoms, cloud infrastructure, search advertising.
  \item[Monopoly] Single supplier, maximum \(L\), subject to regulatory constraint. Rare and regulated;
    a strong network effect (\cref{ch:growth}) can approximate it temporarily.
\end{description}

\Cref{eq:lerner} assumes a static demand curve and no regulatory constraint on pricing. In
practice, high Lerner indices attract antitrust scrutiny (the standard measure in competition
economics). The formula also assigns one elasticity to the whole market; a firm with multiple
segments faces different elasticities --- and thus different optimal prices --- for each.
Long-run equilibrium in monopolistic competition drives \(L\to 0\) as entrants imitate; the
time to equilibrium depends on how defensible the differentiation is (the subject of
\cref{ch:advantage-and-disruption}).

\begin{example}
The SaaS business faces \(\varepsilon=-1.5\) at the current price. Applying \cref{eq:lerner}:
\[
  L = -\frac{1}{-1.5} \approx 0.67.
\]
A Lerner index of \num{0.67} means \qty{67}{\percent} of the \money{50} plan price is mark-up over
marginal cost. Benchmarked against commodity software (near-zero \(L\)) and enterprise software
(often \(L>0.80\)), this is consistent with monopolistic competition with moderate differentiation.
The managerial implication: further differentiation --- integrations, network effects, certification
--- can raise \(|\varepsilon|\) (make demand less elastic) and thus expand the sustainable \(L\)
without price changes. \textcite{dranove7theconomics} survey the empirical range by industry.
\end{example}

\begin{admonition}{Warning}
Confusing the Lerner index with gross margin is a common and consequential error. Gross margin
is \((P-\mathrm{COGS})/P\); the Lerner index is \((P-\mathrm{MC})/P\), where \(\mathrm{MC}\) is
the true marginal cost --- the cost of serving one additional customer at the margin --- which in
SaaS is near zero. A \qty{70}{\percent} gross margin at scale does not imply \(L=0.70\): the two
converge only when there are no fixed costs and COGS captures all variable costs, neither true
at the growth stage nor at scale.
\end{admonition}

In an oligopoly, the Lerner index generalises to the Cournot model: \(L_i = s_i / |\varepsilon|\),
where \(s_i\) is firm \(i\)'s market share. This implies high-share firms have more pricing power
even at the same market elasticity. The Herfindahl-Hirschman Index (\cref{eq:hhi}) in
\cref{ch:competitive-analysis} quantifies concentration; industry-average \(L\) and HHI are
correlated in most empirical studies.

Lerner-index reasoning is directly applied in value-based pricing (\cref{ch:pricing}) and informs
acquisition efficiency decisions (\cref{ch:paid-acquisition-platforms}). \textcite{pindyck2017}
and \textcite{mankiw2020} derive the standard mark-up rules; \textcite{dranove7theconomics}
develops the oligopoly extensions.

\section{Auction Theory and the Ad Auction}\label{sec:auction-theory}
% complexity: 4 -> tabular suffices (discrete 2-bidder case illustrates structure without curve)

You are bidding for attention in a real-time auction that runs in milliseconds. Knowing the
equilibrium bidding strategy --- and why it is optimal to bid your true value --- is the
difference between efficient spend and systematic overpayment. A first-price auction (pay what
you bid) incentivises shading your bid below true value to avoid the winner's curse; the optimal
shade depends on your beliefs about competitors' values. A Vickrey (second-price) auction
removes that incentive: you pay only the second-highest bid, so bidding your true value is a
weakly dominant strategy regardless of what others do. Google and Meta run generalised
second-price (GSP) auctions with a quality weight --- a per-advertiser multiplier that converts
raw bids into effective bids, altering the clearing price and rank.

In the Vickrey second-price auction with \(n\) bidders, truthful bidding is a dominant strategy:
if you win (your true value \(v\) exceeds all rivals' values), you pay the highest rival's value
\(v_{(2)}\) --- less than \(v\), so you earn a positive surplus. If you lose, your bid had no
impact. No deviation from truthful bidding improves your outcome.

The generalised second-price (GSP) auction with quality weight \(Q_i\) ranks bidders by effective
bid \(b_i Q_i\) and charges the winner the amount needed to beat the next-ranked advertiser:
\begin{equation}\label{eq:gsp}
  \mathrm{CPC}_i \;=\; \frac{b_{i+1}\,Q_{i+1}}{Q_i} + \delta,
\end{equation}
where \(b_{i+1}Q_{i+1}\) is the effective AdRank of the bidder ranked just below \(i\) and
\(\delta\) is a small increment (e.g.\ \$0.01). This structure --- already derived operationally
in \cref{eq:ecpc} --- has a key implication: improving quality \(Q_i\) reduces cost per click at
any given rank.

A two-bidder example with values \(v_A=\money{5}\), \(v_B=\money{3}\) and quality scores
\(Q_A=10\), \(Q_B=8\):
\begin{center}\small
\begin{tabular}{lrrrl}
  \toprule
  Bidder & True value & Quality score & Effective bid & Outcome \\
  \midrule
  A      & \money{5}  & 10            & 50            & Wins; pays \(3\times8/10+0.01=\money{2}.41\) \\
  B      & \money{3}  & 8             & 24            & Loses \\
  \bottomrule
\end{tabular}
\end{center}
Bidder A earns surplus \(\money{5}-\money{2}.41=\money{2}.59\) per click. If A reduced quality to
\(Q_A=6\), effective bid falls to 30; B's effective bid of 24 still loses, but A's CPC rises to
\(3\times8/6+0.01=\money{4}.01\), slashing surplus to \money{0}.99. Quality improvement is
unambiguously profitable for the buyer.

\Cref{eq:gsp} assumes the quality score is fixed and known at auction time. In practice both
Google's Quality Score and Meta's estimated action rate (\cref{eq:metavalue}) are updated
continuously; a high-quality creative today may be superseded by a freshly-trained model's lower
prediction tomorrow. GSP is not incentive-compatible in the same clean sense as Vickrey: bidders
can benefit from bid shading in multi-slot environments \parencite{dranove7theconomics}. The
mechanism also assumes independent private values; correlated values (common in brand-keyword
auctions) open space for collusion or tacit coordination.

\begin{example}
The running SaaS business targets the keyword ``project management software'' on Google.
A competitor (B) with AdRank~16 occupies position 2. With Quality Score \(Q_A=8\):
\[
  \mathrm{CPC} = 16/8 + \$0.01 = \money{2.01}.
\]
At \money{120000} monthly budget and CPC of \money{2.01}, the business buys
\(\num{120000}/2.01\approx\num{59700}\) clicks --- well above the learning-phase minimum.
If optimising creative and landing page lifts \(Q_A\) to 10:
\[
  \mathrm{CPC} = 16/10 + \$0.01 = \money{1.61},\quad
  \text{clicks} \approx \num{74500}\;\text{(+25\%)}.
\]
The same budget buys \qty{25}{\percent} more traffic through quality, not spend. This is the
operational form of \cref{eq:gsp} and is consistent with the CPC derivation in \cref{eq:ecpc},
\cref{ch:paid-acquisition-platforms}.
\end{example}

\begin{admonition}{Warning}
Bidding above true value to win position~1 in a GSP auction reduces surplus per click and can
destroy value if the incremental CTR uplift of position~1 over position~2 does not compensate
for the higher CPC. The decision is a marginal analysis: incremental revenue from higher position
versus incremental cost --- exactly the \(\mathrm{MR}=\mathrm{MC}\) condition of
\cref{eq:elasticity}.
\end{admonition}

The revenue equivalence theorem states that under standard conditions (risk-neutral bidders,
independent private values, symmetric distributions) all standard auctions --- first-price,
second-price, English, Dutch --- generate the same expected revenue for the seller. Google and
Meta's use of quality weights is a departure from this symmetric case: it allows the platform to
extract more revenue by ranking higher-quality, lower-bidding advertisers who generate more
user-value (and thus more platform-side surplus) than pure price would imply. The mechanism
design question --- which auction maximises a weighted sum of advertiser and platform welfare ---
is active in the literature \parencite{mankiw2020}.

The GSP structure of \cref{eq:gsp} is the theoretical backbone of the platform mechanics in
\cref{ch:paid-acquisition-platforms}; the quality score is the AdRank component in \cref{eq:adrank}
and the CPC formula in \cref{eq:ecpc}. \textcite{mankiw2020} and \textcite{dranove7theconomics}
cover the auction-theory foundations; the platform implementation is in \cref{ch:paid-acquisition-platforms}.

\section{Transaction Costs and Firm Boundaries}\label{sec:transaction-costs}
% complexity: 3

Should the firm build its own ad-tech stack, manage its own data pipeline, or use the platforms?
Transaction cost economics gives the make-or-buy decision a rigorous frame: internalise when
markets are too costly to use; outsource when they are not. Markets are efficient allocation
mechanisms, but using them is not free. Writing and enforcing contracts, monitoring counterparty
performance, and switching vendors all consume resources. When those transaction costs exceed the
costs of managing an activity inside the firm, internalisation is the efficient choice. Coase's
1937 insight is that firm boundaries are set at the margin where the cost of one more internal
transaction equals the cost of one more market transaction.

The make-or-buy decision is assessed along four dimensions:
\begin{description}
  \item[Asset specificity] How specialised is the asset to this relationship? High specificity
    (custom integrations, proprietary data) creates hold-up risk in a market relationship ---
    the counterparty can extract rents at renegotiation. Internalise.
  \item[Uncertainty] Harder-to-contract environments (rapidly changing tech, unpredictable volumes)
    raise the cost of market contracts. Internalise or use relational (long-term, trust-based) contracting.
  \item[Frequency] Frequent, standardised transactions lower the per-unit cost of market contracting
    relative to the fixed cost of internal governance. Outsource.
  \item[Measurement difficulty] When output quality is hard to observe and verify, market contracts
    produce moral hazard (the counterparty shirks because quality is not punished). Internalise
    or invest in monitoring.
\end{description}
\parencite{coase1937nature}.

The framework assumes that internalisation actually solves the governance problem. Large firms face
internal transaction costs too --- bureaucracy, misaligned incentives, coordination failures ---
that can exceed market costs at scale. The optimal boundary is thus not ``big firm vs.\ market''
but a continuous cost-comparison at each margin of organisational scope. Platform ecosystems create
a third mode: governed markets (App Store, Google Ads API) that combine market pricing with
platform-enforced standards, partially internalising asset-specificity risk without full vertical
integration \parencite{dranove7theconomics}.

\begin{example}
The SaaS business considers whether to build a first-party attribution stack or rely on the
platforms' native conversion APIs. Applying the framework:
\begin{itemize}
  \item \emph{Asset specificity}: first-party conversion data is highly specific (proprietary customer
    identifiers, server-side events); high hold-up risk if stored exclusively in a third-party tool.
    Signal: build.
  \item \emph{Uncertainty}: privacy regulations (cookieless web) are shifting fast; a custom stack
    requires ongoing engineering investment. Signal: buy or use platform API (lower maintenance).
  \item \emph{Frequency}: conversion events happen continuously at scale; standardised APIs handle
    them cheaply. Signal: buy.
  \item \emph{Measurement difficulty}: platform-reported conversions may overcount due to attribution
    windows; verification requires an independent measurement layer. Signal: build the verification
    layer; outsource the delivery.
\end{itemize}
Decision: use the platform Conversions API (\cref{ch:attribution-and-measurement}) for delivery,
build a first-party data layer for verification and modelling. The hybrid minimises both transaction
costs and hold-up risk.
\end{example}

\begin{admonition}{Warning}
Build-vs.-buy decisions made on sunk cost rather than marginal future costs are systematically
biased. A firm that has already invested in a proprietary stack will underestimate ongoing
maintenance cost and overestimate switching cost, both of which favour the status quo regardless
of current economics. Loss aversion (\cref{eq:prospect}) amplifies this: the prospect of
``losing'' the existing investment is coded as a certain loss, weighted at \(\lambda\approx2.25\).
\end{admonition}

Vertical integration and platform governance are extensions of the Coasian framework applied to
digital ecosystems. A platform that owns both the marketplace and a competing product faces the
same make-or-buy problem as a buyer, but reversed: it internalises when it can earn more from
self-supply than from the externality it generates by enabling third-party competition. This is
the economics behind the app-store antitrust disputes: is the platform's vertical integration
(taking a commission, favouring its own products) a welfare-improving reduction in transaction
costs or an anticompetitive foreclosure? \textcite{coase1937nature} provides the foundational
framework; \textcite{dranove7theconomics} apply it to platform markets.

The make-or-buy analysis recurs in operations strategy (\cref{ch:operations}) and in the
discussion of marketing-technology stack decisions in \cref{ch:attribution-and-measurement}.
\textcite{coase1937nature} is the foundational text; \textcite{dranove7theconomics} extend it to
digital markets.
